% ------------------------------------------------------------------
% Section: Support circuits
% Paper:   rotor_window_closures (Erdos #23 series, companion to
%          "Shortest-geodesic supports in triangle-free maximum cuts")
% Drafted from archive sources:
%   problems/23/writeup/WALL_ATTACK_R44_GPTPRO56.md  (section 1)
%   problems/23/writeup/arxiv/shortest_support_obstructions/main.tex
% Assumes the companion preamble: amsthm environments theorem /
% proposition / lemma / corollary / definition / remark numbered by
% section, and the macros \supp, \bip, \mc.
% Bibliography keys used (add at assembly):
%   Ferudun26supports : A. Ferudun, "Shortest-geodesic supports in
%     triangle-free maximum cuts: an infinite Hall obstruction and the
%     first minimal footprints", 2026 (the companion paper).
%   EdmondsFulkerson65 : J. Edmonds and D. R. Fulkerson, "Transversals
%     and matroid partition", J. Res. Nat. Bur. Standards Sect. B 69B
%     (1965), 147--153.
% ------------------------------------------------------------------

\section{Support circuits}\label{sec:support-circuits}

Fix a cut \(V(G)=V_0\mathbin{\dot\cup}V_1\) of a triangle-free graph
\(G\), write \(B\) for the crossing edges and \(M=E(G)\setminus B\) for
the monochromatic edges, and for \(xy\in M\) with \(d_B(x,y)=4\) let
\(\supp_B(xy)\subseteq B\) be the set of edges lying on shortest
\(x\)-\(y\) paths in \(B\), as in the companion paper
\cite{Ferudun26supports}. A failure of the Hall condition
\(|\supp_B(\mathcal S)|\ge|\mathcal S|\) localizes to an
inclusion-minimal violating family, and the Minimal footprint lemma of
\cite{Ferudun26supports} lists the graph-level constraints such a family
must satisfy. In this section we isolate the part of that lemma that
depends on inclusion-minimality alone. It is a statement about abstract
finite set systems; no graph structure enters. Besides recovering the
counting and multiplicity clauses of the footprint lemma, the abstract
form yields two further constraints used repeatedly in later sections:
connectivity of the atom--edge incidence graph, and a family of perfect
systems of distinct representatives obtained by deleting any single
atom. All proofs are elementary and complete.

\subsection{The abstract identity}

Let \(A\) be a finite index set, whose elements we call \emph{atoms},
and let \(\mathcal F=(F_a)_{a\in A}\) be a family of finite sets. For
\(S\subseteq A\) write
\[
 F(S)=\bigcup_{a\in S}F_a,
 \qquad
 F^{*}=F(A).
\]
In the intended application the atoms are bad edges \(xy\in M\) with
\(d_B(x,y)=4\) and \(F_a=\supp_B(a)\), but nothing below uses this.

\begin{definition}\label{def:circuit}
The family \(\mathcal F\) is \emph{deficient} if \(|F^{*}|<|A|\). It is
a \emph{support circuit} if it is deficient and every proper subfamily
satisfies the Hall condition:
\[
 |F(S)|\ge|S|
 \qquad\text{for every } S\subsetneq A.
\]
\end{definition}

Support circuits are exactly the inclusion-minimal deficient families:
a deficient proper subfamily is precisely a violation of the displayed
condition.

\begin{theorem}[Circuit identity]\label{thm:circuit}
Let \(\mathcal F=(F_a)_{a\in A}\) be a support circuit with
\(m=|A|\ge2\). Then:
\begin{enumerate}[label=\textup{(\roman*)}]
\item \(|F^{*}|=m-1\);
\item \(F(A\setminus\{a\})=F^{*}\) for every \(a\in A\);
\item every \(f\in F^{*}\) lies in at least two of the sets \(F_a\);
\item the bipartite incidence graph on \(A\mathbin{\dot\cup}F^{*}\),
with \(a\) adjacent to \(f\) when \(f\in F_a\), is connected;
\item for every \(a\in A\) there is a bijection
\(\varphi_a\colon A\setminus\{a\}\to F^{*}\) with
\(\varphi_a(b)\in F_b\) for all \(b\neq a\).
\end{enumerate}
\end{theorem}

\begin{proof}
(i), (ii). Fix \(a\in A\). The proper subfamily \(A\setminus\{a\}\)
satisfies the Hall condition, so
\[
 m-1\le|F(A\setminus\{a\})|\le|F^{*}|\le m-1,
\]
the last inequality by deficiency and integrality. Equality holds
throughout; in particular \(|F^{*}|=m-1\), and
\(F(A\setminus\{a\})\) is a subset of \(F^{*}\) of full cardinality,
hence equal to it.

(iii). Let \(f\in F^{*}\), say \(f\in F_a\). By (ii),
\(f\in F(A\setminus\{a\})\), so \(f\in F_b\) for some \(b\neq a\).

(iv). Since \(m\ge2\), every singleton \(\{a\}\) is a proper subfamily,
so \(|F_a|\ge1\); and every \(f\in F^{*}\) lies in some \(F_a\). Hence
every connected component of the incidence graph contains at least one
atom. Suppose there are two or more components. Splitting the
components into two nonempty groups partitions \(A\) into nonempty sets
\(A_1,A_2\) with \(F(A_1)\cap F(A_2)=\emptyset\): an element in both
unions would be adjacent to atoms in distinct components. Both \(A_i\)
are proper subfamilies, so
\[
 |F^{*}|=|F(A_1)|+|F(A_2)|\ge|A_1|+|A_2|=m,
\]
contradicting (i).

(v). Every \(S\subseteq A\setminus\{a\}\) is a proper subfamily of
\(A\), so \(|F(S)|\ge|S|\). By Hall's theorem \cite{Hall35} the family
\((F_b)_{b\in A\setminus\{a\}}\) has a system of distinct
representatives: an injection \(\varphi_a\) with
\(\varphi_a(b)\in F_b\). Its image is a subset of \(F^{*}\) of size
\(m-1=|F^{*}|\), so \(\varphi_a\) is a bijection onto \(F^{*}\).
\end{proof}

\begin{remark}\label{rem:intrinsic}
The identity in (i) is intrinsic to the circuit: if a support circuit
\(A\) sits inside a larger ambient family \(\mathcal M\), then
\(|F^{*}|=|A|-1\), not \(|\mathcal M|-1\). All counting in later
sections is against the circuit size \(m\).
\end{remark}

\begin{remark}
The hypothesis \(m\ge2\) only excludes the degenerate circuit
\(A=\{a\}\), \(F_a=\emptyset\), for which (i), (ii) still hold
trivially. In the geodesic application \(m\ge5\)
(Corollary~\ref{cor:sizes}).
\end{remark}

Clause (v) says more than the existence of a matching: after deleting
\emph{any} one atom, the remaining \(m-1\) atoms match onto \emph{all}
of \(F^{*}\), each atom represented inside its own set. This rigidity
is the main tool of the window analyses in later sections. It also
identifies support circuits in classical terms.

\begin{corollary}\label{cor:sdr}
A family \((F_a)_{a\in A}\) is a support circuit if and only if it has
no system of distinct representatives while every proper subfamily has
one.
\end{corollary}

\begin{proof}
By Hall's theorem, a finite family has a system of distinct
representatives if and only if all of its subfamilies satisfy the Hall
condition.

Let \(A\) be a support circuit. Deficiency is a violation of the Hall
condition at \(S=A\), so \(A\) has no system of distinct
representatives. Every subfamily of a proper subfamily of \(A\) is
itself a proper subfamily of \(A\) and hence satisfies the Hall
condition, so every proper subfamily has a system of distinct
representatives.

Conversely, suppose \(A\) has no system of distinct representatives
while every proper subfamily has one. Some \(S\subseteq A\) violates
the Hall condition; \(S\) cannot be proper, since a proper subfamily
has a system of distinct representatives and therefore satisfies
\(|F(S)|\ge|S|\). Hence \(A\) itself is deficient, and every proper
subfamily satisfies the Hall condition.
\end{proof}

\begin{remark}\label{rem:matroid}
The subfamilies of an ambient family \(\mathcal M\) admitting systems
of distinct representatives are the independent sets of a matroid, the
transversal matroid induced by \((F_a)_{a\in\mathcal M}\)
\cite{EdmondsFulkerson65}. Corollary~\ref{cor:sdr} says that the
support circuits inside \(\mathcal M\) are exactly the circuits of this
matroid, and Theorem~\ref{thm:circuit}(v) computes the rank of a
circuit deleted by one element: the deletion matchings are perfect onto
the ground union \(F^{*}\).
\end{remark}

\subsection{The graph form}

Specializing to shortest-geodesic supports recovers the Minimal
footprint lemma of \cite{Ferudun26supports} and extends it by the two
new clauses. We restate the full result in the paper's setting; the
genuinely graph-level clauses (c), (d) have the same short proofs as in
\cite{Ferudun26supports} and are included for completeness.

\begin{corollary}[Graph support circuits]\label{cor:graph}
Let \(G\) be triangle-free with a fixed cut, let \(B\) be connected,
and let \(\mathcal A\subseteq\{xy\in M:d_B(x,y)=4\}\) be
inclusion-minimal with
\[
 |\supp_B(\mathcal A)|<|\mathcal A|=m.
\]
Write \(F=\supp_B(\mathcal A)\), viewed as the subgraph it spans. Then:
\begin{enumerate}[label=\textup{(\alph*)}]
\item \(|E(F)|=m-1\), and \(F\) is connected and bipartite;
\item every edge of \(F\) lies in at least two atom supports, and
\(\supp_B(\mathcal A\setminus\{a\})=E(F)\) for every
\(a\in\mathcal A\);
\item every \(xy\in\mathcal A\) satisfies \(d_F(x,y)=4\), and
\(\supp_B(xy)\) is the union of all length-four \(x\)-\(y\) paths in
\(F\);
\item the graph \((V(F),\mathcal A)\) is triangle-free;
\item for every \(a\in\mathcal A\) there is a bijection from
\(\mathcal A\setminus\{a\}\) onto \(E(F)\) assigning to each atom an
edge of its own support.
\end{enumerate}
\end{corollary}

\begin{proof}
The family \((\supp_B(a))_{a\in\mathcal A}\) is a support circuit, and
every support is nonempty because \(d_B(x,y)=4\), so \(m\ge2\) and
Theorem~\ref{thm:circuit} applies with \(F^{*}=E(F)\).

(a) The count is Theorem~\ref{thm:circuit}(i), and \(F\subseteq B\) is
bipartite. For connectivity: each support \(\supp_B(xy)\) is a
connected edge set, being a union of \(x\)-\(y\) paths that all contain
\(x\). If two atoms are adjacent to a common edge \(f\) in the
incidence graph, their supports share the edge \(f\) and hence their
union is connected. Walking along paths of the incidence graph, which
is connected by Theorem~\ref{thm:circuit}(iv), connects all supports,
and every edge of \(F\) lies in some support.

(b) Theorem~\ref{thm:circuit}(iii) and (ii).

(c) The edges of any shortest \(x\)-\(y\) path in \(B\) lie in
\(\supp_B(xy)\subseteq E(F)\), so \(d_F(x,y)\le4\). Since
\(F\subseteq B\) and \(x,y\) lie on the same shore, every \(x\)-\(y\)
walk in \(F\) has even length, and \(d_F(x,y)=2\) would give a common
neighbour forming a triangle with the edge \(xy\). Hence
\(d_F(x,y)=4\). A length-four \(x\)-\(y\) path in \(F\) is a
length-four \(x\)-\(y\) path in \(B\), hence shortest, hence contained
in \(\supp_B(xy)\); conversely every shortest \(x\)-\(y\) path in \(B\)
lies in \(\supp_B(xy)\subseteq E(F)\) and has length
\(4=d_F(x,y)\), so it is one of the length-four \(F\)-paths.

(d) \(\mathcal A\subseteq E(G)\) and \(G\) is triangle-free; the
endpoints of atoms lie in \(V(F)\) by (c).

(e) Theorem~\ref{thm:circuit}(v).
\end{proof}

\begin{corollary}[First sizes]\label{cor:sizes}
In the setting of Corollary~\textup{\ref{cor:graph}}, \(m\ge5\), and
the average multiplicity of a support edge satisfies
\[
 \frac{1}{|E(F)|}\sum_{a\in\mathcal A}|\supp_B(a)|
 \;\ge\;\frac{4m}{m-1}\;>\;4.
\]
\end{corollary}

\begin{proof}
Each support contains the four edges of a shortest path, so
\(4\le|\supp_B(a)|\le|E(F)|=m-1\), giving \(m\ge5\); and
\(\sum_a|\supp_B(a)|\ge4m\) with \(|E(F)|=m-1\).
\end{proof}

\begin{remark}
Corollary~\ref{cor:graph}(a)--(d) recovers the Minimal footprint lemma
of \cite{Ferudun26supports}, where it feeds the exact classification of
the footprints with \(m\le10\); clause (e), the incidence connectivity
behind (a), and the deletion-union identity in (b) (which occurs in the
proof, but not the statement, of the companion lemma) are the
additional structure carried by the circuit form. Sharper counts are
derived in later sections for \emph{selected} supports, in which each
atom retains a single geodesic; these depend on
the owner--profile apparatus. The
results of this section constrain the geometry of a minimal Hall
violation; they give no upper bound on \(\bip(G)\).
\end{remark}
